\chapter{While traveling}

In this section we're going to cover all the things you
have to pay attention to or take care of while on the
move. Is it necessary to connect to that conference's
Wi-Fi? Have you checked if there is a camera that might be
recording your keystrokes? Leaving your computer
unattended or just running whatever software your
hackathon teammate asks you to download in order for your
promising project to win might not be the best idea, it
could result in compromise.

\section{Network safety}

Treat every network as potentially hostile. Whenever
possible, use a cellular connection or a personal hotspot
instead of public Wi-Fi -- your mobile data is encrypted
and safer than an open café or hotel network. If you must
use public Wi-Fi (hotel, airport, conference), verify the
network name with staff and disable auto-connect
features\footnote{\url{https://www.cisa.gov/sites/default/files/publications/Cybersecurity\%20While\%20Traveling.pdf}}
so your devices don't join networks without prompting you.

Turn off Wi-Fi and Bluetooth on your phone and laptop when
you're not using them; this reduces the risk of
unsolicited connections or Bluetooth-based attacks. Also,
disable any device-to-device sharing like AirDrop or
Nearby Share to prevent strangers from sending you
files.\footnote{\url{https://www.cypherock.com/blogs/post-safe-vacation-tips-while-traveling-with-crypto}}

Use a trustworthy VPN for an extra layer of encryption for
your internet traffic, although in most cases by using an
updated device, safe hardcoded DNS records, and ensuring
SSL while browsing (enforce HTTPS-only-modes or adding
``http://*'' to your uBlock list) might be enough. A
reputable, non-logging VPN (or your company's VPN) helps
protect against snooping on public networks, especially if
you're handling highly sensitive work and using several
communication channels.

Using a personal portable router combined with a trusted
VPN adds a strong layer of security when connecting to
public Wi-Fi networks. This setup creates a private,
encrypted tunnel between your devices and the internet,
minimizing exposure to malicious actors on shared
networks. Whenever possible, prefer mobile data over
Wi-Fi, as cellular networks provide better encryption and
isolation by default. If you must use Wi-Fi, disable
automatic connections and ensure you connect only to
verified, trusted networks to reduce risk.

\section{Device handling}

Keep them close and protected. Never leave your devices
(laptops, phones, hardware wallets) unattended or unlocked
in public. Keep them on your person or within sight
whenever possible. In a conference or cafe, use a cable
lock for your laptop if you must step away briefly, take
it with you or get someone you trust to watch it over for
you.

In hotels/Airbnbs, use the room safe for small devices
when you're out. These safes can easily be opened by an
experienced thief or rogue/malicious staff. Consider a
portable travel safe/bag you can lock to a fixed object.

A portable door
lock\footnote{\url{https://www.travelandleisure.com/sabre-portable-door-lock-review-7643179}}
or door jammer for your room can add an extra barrier
against intruders (useful in rentals without chain locks).
This simple gadget prevents anyone (even with a key) from
opening your door while you're inside. When out and about,
be mindful of pickpockets -- use bags that zip and
consider a subtle anti-theft backpack for expensive
gadgets or important assets.

For hardware wallets or Yubikeys, a good practice is to
separate them from the device they authenticate: e.g.\
don't carry your Yubikey on the same keychain as your
laptop bag key; keep it hidden on you. And, of course,
keep devices powered off or locked when not in use --
enable short auto-lock
timeouts\footnote{\url{https://www.cisa.gov/sites/default/files/publications/Cybersecurity\%20While\%20Traveling.pdf}}
on your phone/laptop so they aren't unlocked if snatched.
You can also do this before leaving as well.

\section{Beware of unknown USB charging}

Avoid plug-and-play charging stations at airports or
malls. The risk of ``juice
jacking'',\footnote{\url{https://www.cypherock.com/blogs/post-safe-vacation-tips-while-traveling-with-crypto}}
although small, is that a malicious charging kiosk or
cable can inject malware or siphon data when you connect
via USB. Depending on your profile risk, you might
encounter technologies that mimic different kind of cables
and accessories with the same purpose. These exist out of
the box, and have been widely used by red-teamers and
pentesters (re OMG cable).

Stick to using your own charger plugged into a power
outlet, or use a USB data blocker (a little adapter that
only passes power, not data). Similarly, do not plug
unknown USB drives into your laptop -- if someone hands
you a free USB stick at an event, assume it could be a
trap.

USBGuard\footnote{\url{https://github.com/USBGuard/usbguard}}
software (for Linux) or equivalent can be used to restrict
USB device access on your computer, allowing only
whitelisted devices. This tool can prevent an unknown USB
from automatically tampering with your system by requiring
authorization for new devices. At a minimum, disable any
``USB autorun'' features and consider locking down ports
if your OS allows.

\section{Screen privacy and social engineering}

Practice situational awareness when working in public
spaces. Shoulder
surfing\footnote{\url{https://www.trio.so/blog/shoulder-surfing-in-computer-security}}
is a real threat -- someone nearby could be watching you
enter passwords or PIN codes.

Use a privacy screen filter on your laptop or phone to
narrow the viewing angle. This makes it much harder for
anyone not directly in front of your screen to read it.
Sit with your back to a wall when possible, and shield the
keypad with your body or hand when typing sensitive info.

In crowded conferences or airports, be cautious if someone
you don't know strikes up conversation -- might be trying
to distract you or talk about crypto; scammers might
engage you to glean info or even attempt to get you to
unlock your device. Don't log in to critical accounts in
front of others -- you never know who's looking.

Also be mindful of phishing attempts: traveling users are
prime targets for fake ``urgent'' emails or messages.
Double-check any unusual prompts before entering
credentials, especially if you're on untrusted
Wi-Fi\footnote{\url{https://www.cisa.gov/sites/default/files/publications/Cybersecurity\%20While\%20Traveling.pdf}}
(use your VPN and look for HTTPS).

\section{Maintain OpSec in public}

While traveling, blend in and stay discreet. Refrain from
discussing sensitive matters in public areas where
eavesdropping is possible, or directly sharing things like
where you are staying to people you don't know. Even hotel
lobbies or ride-shares might not be secure for private
discussions.

When meeting people, don't give your phone to others to
type down their socials, and remember to disable default
options like Telegram's sharing phone number when adding a
contact.

Remember that hidden cameras or microphones could exist in
unfamiliar environments. It's rare but not unheard of,
especially in Airbnbs or rented
spaces\footnote{\url{https://www.washingtonpost.com/travel/2024/05/22/hidden-cameras-airnbnb-rental-properties}}
-- malicious hosts have hidden cameras in items like smoke
detectors, clock radios, or USB chargers. Give your
accommodations a scan: look for odd or extra devices
plugged in (especially facing beds or desks) and cover or
unplug them if suspicious. You can also play ambient noise
(or use a noise generator app) during confidential
conversations to thwart any listening devices.

\subsection{Keep a low profile}

As mentioned, don't flaunt crypto wealth or
gear.\footnote{\url{https://www.unchained.com/blog/7-tips-traveling-bitcoin}}
For example, if attending a blockchain conference,
consider using an alias on your name badge that doesn't
explicitly say your company or title, and don't display
that badge outside the venue. When moving around, secure
your laptop in a nondescript sleeve or bag (instead of one
with a well-known conference brand). The goal is to avoid
drawing the attention of both petty thieves and more
organized attackers by limiting the signals that you're a
high-value crypto target.

Don't rely on security through obscurity. Going
``stealth'' doesn't make you immune to attacks. The
suggestions in this section are meant to complement, not
replace, the other security measures.

\section{Traveling with high-value crypto or duties}

If you must make crypto transactions or access sensitive
systems while on the road, do so with caution. Use trusted
hardware and networks: e.g.\ if you need to send a
transaction, use your hardware wallet on your own laptop
(never a shared computer), on a secured connection.

Be aware of surveillance at events -- attackers have been
known to watch for people handling sensitive info. If you
need to access a seed or enter a recovery phrase, do it in
a private, secure setting (never over public Wi-Fi or in
view of anyone, including cameras).

Consider that everyone knows you own crypto at a crypto
event,\footnote{\url{https://www.unchained.com/blog/7-tips-traveling-bitcoin}}
so your threat profile is elevated. Adjust your security:
for instance, enable a passphrase or a pin on any
single-signature wallet you carry so that even if someone
obtains your hardware wallet, they can't access funds
without that passphrase. For large amounts, rely on
multi-sigs, verifying payloads before signing -- you might
carry one key on you and leave another key(s) with trusted
parties so no single person has all signing power while
traveling. In short, treat any on-trip crypto operations
with more care than you would in the office or at home.

\section{While presenting / doing public appearances}

One often overlooked risk is the exposure caused by
presenting or hosting technical workshops in public.
Without properly hardening or isolating your computer
before setting up, you may unintentionally expose network
services to hostile environments or reveal sensitive
information on-screen. Always prepare a clean, minimal
environment and verify no confidential data or open ports
are accessible before connecting to unfamiliar networks or
projecting your screen.

\section{Physical safety and common sense}

Operational security also has a physical aspect. Trust
your instincts and normal travel safety rules: stick to
well-lit and populated areas if carrying devices at night,
don't let strangers ``shoulder surf'' your ATM or credit
card PIN (check for skimmers, fake interfaces), and keep
your travel documents secure since identity theft can be
as damaging as device theft.

When possible, store your devices in secure lockers (many
hotels or conferences offer charging lockers). Never leave
them in open areas.

Beware of the classic ``evil maid'' scenario in hotels
(where someone might tamper with your laptop in your
room): using tamper-evident tape or seals on your laptop
case can help detect this, though it's mostly a concern
for high-risk targets. If you have tamper-evident stickers
or tamper-evident bags, you can seal your device in them
overnight -- any attempt to open or remove the device will
leave a visible trace. While not foolproof against a
determined adversary, it raises the bar and can deter
casual snooping.

Petty thieves may look beyond obvious valuables. Simply
locking items in a dorm safe or hiding them at home might
deter casual theft, but savvy criminals often search
inside books, behind electrical outlets, or within
patterns on walls or furniture to find hidden stashes.
Consider unconventional hiding spots and avoid predictable
storage methods. Layer your physical security measures
with tamper-evident seals or discreet decoy containers to
raise the effort required for unauthorized access.

Above all, maintain an alert posture: be aware of who's
around when you're working, and if something feels off
(like someone persistently hovering or a device acting
strangely), don't ignore it. You can always relocate,
power down your device, or otherwise cut off exposure at
the first sign of potential trouble.

There's more information for high profile individuals in
the last chapter, page~\pageref{ch:high-profile}.
